% =====================================================================
% CAP. 11 - NIVEL 4 - PROJETO, DIAGNOSTICO E CIRCUITOS INTEGRADOS
% =====================================================================
\blocodesafio

\begin{questao}[4]{}
Um sensor pode ser modelado por uma fonte de Thévenin de $\SI{5}{\volt}$ com
resistência interna $R_s=\SI{1}{\kilo\ohm}$. Ao fechar a chave em $t=0$, ele
alimenta uma entrada com $R_L=\SI{9}{\kilo\ohm}$ em paralelo com um capacitor
$C$, inicialmente descarregado. O sistema digital reconhece nível alto quando
$v_C\ge\SI{3}{\volt}$ e exige que isso ocorra em no máximo
$\SI{20}{\milli\second}$.

\circuitoqq{%
\begin{circuitikz}[scale=0.90,transform shape,line width=0.8pt,every node/.append style={font=\scriptsize}]
\draw (0,0) to[american voltage source,l_={$\SI{5}{\volt}$}] (0,3)
      to[R,l={$R_s=\SI{1}{\kilo\ohm}$}] (2.5,3)
      to[switch] (4.3,3) -- (6.8,3) coordinate(a);
\node[anchor=south] at (3.4,3.30) {fecha em $t=0$};
\draw (a) to[R,l_={$R_L=\SI{9}{\kilo\ohm}$}] (6.8,0) -- (0,0);
\draw (a) -- (9.2,3) to[C,v^>={$v_C$}] (9.2,0) -- (6.8,0);
\node[anchor=east] at (8.95,1.5) {$C$};
\end{circuitikz}}{Entrada RC carregada por uma fonte com resistência interna e carga finita.}

\begin{itensq}
\item Determine o valor final de $v_C$ e a resistência de Thévenin vista pelo capacitor.
\item Calcule o maior valor de $C$ que atende ao tempo especificado.
\item Escolha entre $\SI{18}{\micro\farad}$ e $\SI{22}{\micro\farad}$ e verifique o tempo de cruzamento de \SI{3}{\volt}.
\item Determine a corrente fornecida pela fonte em $0^+$ e em regime permanente.
\end{itensq}

\resposta{$V_f=\SI{4.5}{\volt}$; $R_{\mathrm{Th}}=\SI{900}{\ohm}$;
$C_{\max}\approx\SI{20.2}{\micro\farad}$; escolha: \SI{18}{\micro\farad};
$t_{3V}\approx\SI{17.8}{\milli\second}$;
$i_s(0^+)=\SI{5}{\milli\ampere}$; $i_s(\infty)=\SI{0.5}{\milli\ampere}$}
\resolucao{Em regime permanente o capacitor é aberto, então
\[
V_f=5\frac{9}{1+9}=\SI{4.5}{\volt}.
\]
Suprimindo a fonte, a resistência vista pelo capacitor é
$R_{\mathrm{Th}}=R_s\parallel R_L=1\parallel9=\SI{0.9}{\kilo\ohm}$.
Como o capacitor parte descarregado,
$v_C(t)=4{,}5(1-e^{-t/(R_{\mathrm{Th}}C)})$. A condição limite é
$3=4{,}5(1-e^{-0{,}02/(900C)})$, ou seja $e^{-0{,}02/(900C)}=1/3$. Logo
\[
C_{\max}=\frac{0{,}02}{900\ln3}\approx\SI{20.2}{\micro\farad}.
\]
O valor de \SI{22}{\micro\farad} viola a especificação; com
$C=\SI{18}{\micro\farad}$, $\tau=\SI{16.2}{\milli\second}$ e
$t_{3V}=\tau\ln3\approx\SI{17.8}{\milli\second}$. Em $0^+$ o capacitor se
comporta como curto, logo $i_s=5/1\,\mathrm{k\Omega}=\SI{5}{\milli\ampere}$.
Em regime, a corrente é $5/(1+9)\,\mathrm{k\Omega}=\SI{0.5}{\milli\ampere}$.}
\end{questao}

\begin{questao}[4]{}
Um barramento CC possui um capacitor de $\SI{470}{\micro\farad}$ carregado a
$\SI{325}{\volt}$. Quando a alimentação é desconectada, um resistor de descarga
$R_b$ deve reduzir a tensão para no máximo $\SI{50}{\volt}$ em
$\SI{60}{\second}$. Despreze outras cargas.

\circuitoqq{%
\begin{circuitikz}[scale=0.94,transform shape,line width=0.8pt,every node/.append style={font=\scriptsize}]
\draw (0,0) to[american voltage source,l_={$\SI{325}{\volt}$}] (0,3)
      to[switch] (2.0,3) -- (5.2,3) coordinate(a);
\node[anchor=south] at (1.0,3.30) {desliga em $t=0$};
\draw (a) to[C,v^>={$v_C$}] (5.2,0) -- (0,0);
\node[anchor=east] at (4.95,1.5) {$C=\SI{470}{\micro\farad}$};
\draw (a) -- (7.8,3) to[R,l_={$R_b$}] (7.8,0) -- (5.2,0);
\end{circuitikz}}{Resistor de descarga de um barramento capacitivo.}

\begin{itensq}
\item Determine o maior valor permitido de $R_b$.
\item Escolha um valor comercial de $\SI{68}{\kilo\ohm}$ e verifique o tempo até \SI{50}{\volt}.
\item Calcule a potência instantânea inicial no resistor e a energia total que ele dissipará.
\end{itensq}

\resposta{$R_{b,\max}\approx\SI{68.2}{\kilo\ohm}$; com \SI{68}{\kilo\ohm},
$t_{50V}\approx\SI{59.8}{\second}$; $p_R(0^+)\approx\SI{1.55}{\watt}$;
$E\approx\SI{24.8}{\joule}$}
\resolucao{Na descarga, $v_C(t)=V_0e^{-t/(R_bC)}$. Impor
$50=325e^{-60/(R_bC)}$ fornece
\[
R_b=\frac{-60}{C\ln(50/325)}\approx\SI{68.2}{\kilo\ohm}.
\]
Assim, \SI{68}{\kilo\ohm} atende à especificação. Com esse valor,
$t=R_bC\ln(325/50)\approx\SI{59.8}{\second}$. A potência inicial é
$p_R(0^+)=V_0^2/R_b\approx\SI{1.55}{\watt}$. Toda a energia inicialmente
armazenada no capacitor é dissipada no resistor:
\[
E=\frac12CV_0^2
=\frac12(470\times10^{-6})(325)^2\approx\SI{24.8}{\joule}.
\]
O projeto real deve ainda considerar a tensão máxima do resistor e sua capacidade
de suportar a energia do pulso de descarga.}
\end{questao}

\begin{questao}[4]{}
Projete o RLC em série abaixo para que a resposta do capacitor a um degrau seja
subamortecida com $\zeta=0{,}5$ e tempo de acomodação de \SI{2}{\percent} não
superior a $\SI{40}{\milli\second}$. Adote
$C=\SI{100}{\micro\farad}$ e use a aproximação
$t_s\approx4/(\zeta\omega_0)$. Escolha o menor $\omega_0$ que atende à
especificação. Determine $L$, $R$, o sobressinal percentual e o instante do
primeiro pico.

\circuitoqq{%
\begin{circuitikz}[scale=0.90,transform shape,line width=0.8pt,every node/.append style={font=\scriptsize}]
\draw (0,0) to[american voltage source,l_={$V_su(t)$}] (0,3)
      to[switch] (1.8,3)
      to[R,l={$R$}] (4.0,3)
      to[L,l={$L$}] (6.5,3)
      to[C,v^>={$v_C$}] (9.0,3)
      -- (9.0,0) -- (0,0);
\node[anchor=south] at (0.9,3.30) {fecha em $t=0$};
\node[anchor=south] at (7.75,3.30) {$C=\SI{100}{\micro\farad}$};
\end{circuitikz}}{Projeto de um circuito RLC em série a partir de especificações transitórias.}

\resposta{$\omega_0=\SI{200}{\radian\per\second}$;
$L=\SI{0.25}{\henry}$; $R=\SI{50}{\ohm}$;
$M_p\approx\SI{16.3}{\percent}$; $t_p\approx\SI{18.1}{\milli\second}$}
\resolucao{Usando a especificação no limite,
\[
0{,}040=\frac{4}{0{,}5\omega_0}
\quad\Rightarrow\quad
\omega_0=\SI{200}{\radian\per\second}.
\]
Como $\omega_0=1/\sqrt{LC}$,
$L=1/(C\omega_0^2)=1/(100\times10^{-6}\cdot200^2)=\SI{0.25}{\henry}$.
Para RLC em série, $\zeta=R/(2L\omega_0)$; portanto
$R=2\zeta L\omega_0=\SI{50}{\ohm}$. O sobressinal é
\[
M_p=e^{-\pi\zeta/\sqrt{1-\zeta^2}}\approx0{,}163,
\]
ou \SI{16.3}{\percent}. A frequência amortecida é
$\omega_d=\omega_0\sqrt{1-\zeta^2}\approx\SI{173.2}{\radian\per\second}$ e
$t_p=\pi/\omega_d\approx\SI{18.1}{\milli\second}$.}
\end{questao}

\begin{questao}[4]{}
O circuito RC em escada abaixo possui \textbf{dois capacitores independentes} e,
portanto, é de segunda ordem. Para $v_s(t)=10u(t)\,\si{\volt}$ e condições
iniciais nulas, obtenha as equações de estado em $v_1$ e $v_2$, determine os polos
naturais e encontre $v_2(t)$.

\circuitoqq{%
\begin{circuitikz}[scale=0.94,transform shape,line width=0.8pt,every node/.append style={font=\scriptsize}]
\draw (0,0) to[american voltage source] (0,3)
      to[R,l={$R_1=\SI{1}{\kilo\ohm}$}] (3.2,3) coordinate(a)
      to[R,l={$R_2=\SI{2}{\kilo\ohm}$}] (6.5,3) coordinate(b);
\node[anchor=east] at (-0.25,1.5) {$10u(t)\,\si{\volt}$};
\draw (a) to[C,v^>={$v_1$}] (3.2,0);
\node[anchor=east] at (2.95,1.5) {$C_1=\SI{100}{\micro\farad}$};
\draw (b) to[C,v^>={$v_2$}] (6.5,0) -- (0,0);
\node[anchor=east] at (6.25,1.5) {$C_2=\SI{50}{\micro\farad}$};
\end{circuitikz}}{Rede RC de segunda ordem com dois elementos armazenadores do mesmo tipo.}

\resposta{$\dot v_1=100u(t)-15v_1+5v_2$;
$\dot v_2=10v_1-10v_2$; polos $-5$ e $-20\ \si{\per\second}$;
$v_2(t)=10-\frac{40}{3}e^{-5t}+\frac{10}{3}e^{-20t}\ \si{\volt}$}
\resolucao{Aplicando KCL aos dois nós,
\[
C_1\dot v_1=\frac{v_s-v_1}{R_1}-\frac{v_1-v_2}{R_2},
\qquad
C_2\dot v_2=\frac{v_1-v_2}{R_2}.
\]
Substituindo os valores,
\[
\dot v_1=100u(t)-15v_1+5v_2,\qquad
\dot v_2=10v_1-10v_2.
\]
A matriz natural é
$A=\begin{bmatrix}-15&5\\10&-10\end{bmatrix}$, cuja equação característica é
$s^2+25s+100=0$, fornecendo $s_1=-5$ e $s_2=-20$. Em regime final,
$v_1(\infty)=v_2(\infty)=\SI{10}{\volt}$. Assim
$v_2=10+Ae^{-5t}+Be^{-20t}$. Das condições $v_2(0)=0$ e
$\dot v_2(0)=0$ resultam $A=-40/3$ e $B=10/3$.}
\end{questao}

\begin{questao}[4]{}
No circuito RLC em paralelo, a fonte de corrente vale $\SI{2}{\ampere}$ por muito
tempo e muda instantaneamente para $\SI{1}{\ampere}$ em $t=0$. Determine a
resposta completa $v(t)$, classifique o amortecimento e encontre o primeiro pico
negativo da tensão.

\circuitoqq{%
\begin{circuitikz}[scale=0.94,transform shape,line width=0.8pt,every node/.append style={font=\scriptsize}]
\draw (0,0) to[american current source] (0,3) -- (7.5,3);
\node[anchor=east,align=right] at (-0.25,1.5) {$I_s$\\$2\,\mathrm{A}\to1\,\mathrm{A}$};
\draw (2.3,3) to[R] (2.3,0);
\node[anchor=east] at (2.05,1.5) {$R=\SI{20}{\ohm}$};
\draw (4.9,3) to[C,v^>={$v$}] (4.9,0);
\node[anchor=east] at (4.65,1.5) {$C=\SI{500}{\micro\farad}$};
\draw (7.5,3) to[L,i>^={$i_L$}] (7.5,0) -- (0,0);
\node[anchor=east] at (7.25,1.5) {$L=\SI{0.2}{\henry}$};
\end{circuitikz}}{RLC em paralelo excitado por uma mudança de nível da fonte de corrente.}

\resposta{Subamortecida;
$v(t)\approx-23{,}09e^{-50t}\sin(86{,}60t)\ \si{\volt}$;
$t_p\approx\SI{12.1}{\milli\second}$;
$v_{\min}\approx-\SI{10.93}{\volt}$}
\resolucao{Antes da mudança, em regime CC, o indutor é curto e conduz toda a
corrente da fonte: $i_L(0^-)=\SI{2}{\ampere}$ e $v(0^-)=0$. Logo
$i_L(0^+)=\SI{2}{\ampere}$ e $v(0^+)=0$. Para $t>0$, o valor final é novamente
$v(\infty)=0$, com $i_L(\infty)=\SI{1}{\ampere}$. Os parâmetros são
\[
\alpha=\frac{1}{2RC}=\SI{50}{\per\second},\qquad
\omega_0=\frac{1}{\sqrt{LC}}=\SI{100}{\radian\per\second},
\]
portanto $\omega_d=\sqrt{100^2-50^2}\approx\SI{86.60}{\radian\per\second}$.
Aplicando KCL em $0^+$,
\[
1=\frac{0}{20}+C\dot v(0^+)+2
\Rightarrow
\dot v(0^+)=-\frac{1}{500\,\mathrm{\mu F}}
=-\SI{2000}{\volt\per\second}.
\]
Como $v(0)=0$, resulta
$v(t)=B e^{-50t}\sin(86{,}60t)$, com
$B=-2000/86{,}60\approx-23{,}09$. O primeiro extremo satisfaz
$\tan(\omega_dt)=\omega_d/\alpha=\sqrt3$, logo
$\omega_dt=\pi/3$ e $t_p\approx\SI{12.1}{\milli\second}$. Substituindo na
resposta, $v_{\min}\approx-\SI{10.93}{\volt}$.}
\end{questao}

\gabarito[Nível 4 --- projeto, diagnóstico e síntese]
